\documentclass{article}
\usepackage{amsmath}
\begin{document}

\begin{abstract}
Reliable manuscript construction requires more than grammatical correction.
This fixture represents a compact journal structure in which the chapter
sequence, local argument, notation, equations, model description, and
experimental workflow are checked after every substantive edit. The method
uses a section-role matrix to keep each section tied to one scientific
question and one reader-facing result, while a notation ledger registers every
mathematical object before use. A subsection audit then records the logical
role of adjacent sentences, the causal chain from the motivating limitation to
the proposed mechanism and its consequence, and the assumptions and dimensions
needed by each formula. It also verifies that model inputs, outputs, states,
mappings, and update order are complete. Training, validation, testing, and
deployment are described as separate flows so that model selection and
preprocessing cannot leak test information. The resulting workflow provides a
repeatable record of revisions and blocks delivery whenever an unresolved
conflict remains. This record also supports transparent review.
\end{abstract}

\section{Introduction}

The main contributions are:

\begin{enumerate}
\item A compact chapter-control procedure aligns each central problem with one
principal contribution and one explicit body result.
\item A seven-part subsection audit records logical, mathematical, model, and
workflow evidence after every substantive edit.
\end{enumerate}

\section{Preliminaries and Problem Formulation}

\subsection{Problem Formulation}

\begin{enumerate}
\item Maintain a scientifically ordered chapter sequence without duplicated
section responsibilities.
\item Verify every revised subsection through the complete writing-quality
loop before advancing.
\end{enumerate}

\section{Proposed Method}

Let $x$ denote the scalar model state and let $y$ denote the scalar output. The
complete fixture model is

\begin{equation}
  y=x.
  \label{eq:writing-loop-fixture}
\end{equation}

The training split estimates the model, the validation split selects the final
configuration, and the untouched test split evaluates the selected
configuration.

\section{Experiments}

The fixture audit checks the recorded split roles and does not report an
unobserved numerical result.

\section{Conclusion}

The chapter and subsection records provide the evidence required by the
writing loop.

\end{document}
